% © 2026 Ing. David Jaroš — CC BY-NC-ND 4.0
%
% This work is licensed under a Creative Commons Attribution-NonCommercial-NoDerivatives
% 4.0 International License (CC BY-NC-ND 4.0).
%
% License History: Earlier drafts (up to v0.3) were released under CC BY 4.0.
% From v0.4 onward, all material is released under CC BY-NC-ND 4.0 to protect
% the integrity of the theoretical work during ongoing academic development.
%
% Particle Spectrum from Theta Excitation Types — UBT Particle Dictionary
% Status: MAP document — every assignment carries a status label (Proved/Conjecture/Open/Speculative)
% Version: 1.0
% Date: 2026-03-05

\documentclass[12pt,a4paper]{article}

\usepackage{amsmath,amssymb,amsthm}
\usepackage{geometry}
\usepackage{booktabs}
\usepackage{array}
\usepackage{hyperref}
\usepackage{enumitem}
\usepackage{xcolor}
\geometry{margin=2.5cm}

% Theorem environments
\newtheorem{theorem}{Theorem}[section]
\newtheorem{lemma}[theorem]{Lemma}
\newtheorem{proposition}[theorem]{Proposition}
\theoremstyle{remark}
\newtheorem{remark}[theorem]{Remark}
\newtheorem{conjecture}[theorem]{Conjecture}

% Notation shorthands
\newcommand{\B}{\mathcal{B}}
\newcommand{\C}{\mathbb{C}}
\newcommand{\R}{\mathbb{R}}
\newcommand{\H}{\mathbb{H}}
\newcommand{\Z}{\mathbb{Z}}
\newcommand{\Tfield}{\Theta}
\newcommand{\re}{\mathrm{Re}}
\newcommand{\im}{\mathrm{Im}}
\newcommand{\Tr}{\mathrm{Tr}}

% Status labels
\newcommand{\proved}{\textbf{[Proved]}}
\newcommand{\conjecture}{\textbf{[Conjecture]}}
\newcommand{\openq}{\textbf{[Open]}}
\newcommand{\speculative}{\textbf{[Speculative]}}

\title{%
  Particle Spectrum from $\Theta$ Excitation Types:\\
  A Map of UBT Particle Assignments\\[0.5ex]
  \large UBT Core Document --- Particle Dictionary v1.0
}

\author{David Jaro\v{s}}

\date{2026-03-05}

\begin{document}

\maketitle

\begin{abstract}
In the Unified Biquaternion Theory (UBT) there is ONE fundamental dynamical
object: the biquaternionic master field
$\Theta(x,\tau)\in\B = \C\otimes_\R\H$, with complex time
$\tau = t + i\psi$.  Every observed particle is identified with a distinct
\emph{type of excitation} of this single field.  This document is a
\textbf{map}, not a proof.  It catalogues four excitation types (A--D),
assigns each known particle type to one or more excitation types, labels
every assignment with its current proof status (\proved, \conjecture,
\openq, or \speculative), counts the degrees of freedom, and states the
open problems honestly.  The goal is to turn conjectures into falsifiable
targets.
\end{abstract}

\tableofcontents

%=============================================================
\section{Introduction: One Field, All Particles}
\label{sec:intro}
%=============================================================

In the Standard Model (SM) of particle physics each particle species comes
with its own independent field: an electron field $\psi_e$, a photon field
$A_\mu$, quark fields $q$, and so forth.  Unification in the SM sense means
finding a gauge group large enough to contain all of them; it does not reduce
their number.

UBT proposes a sharper unification:

\begin{quote}
\emph{There is one field $\Theta(x,\tau)\in\B$.  All particles are different
excitation modes of this field.  Their apparent diversity reflects the
geometric richness of $\B$ over complex time $\tau = t+i\psi$.}
\end{quote}

The decomposition
\begin{equation}
  \Theta = \Theta_{\text{scalar}} + \Theta_{\text{vector}}
           + \Theta_{\text{spinor}} + \Theta_{\text{gauge}}
  \label{eq:decomp}
\end{equation}
(from \texttt{THETA\_FIELD\_DEFINITION.md}) provides the first, coarsest
classification.  This document refines it into four excitation types and
builds a full particle dictionary.

The document is structured as follows: Section~\ref{sec:types} defines the
four excitation types A--D.  Section~\ref{sec:dict} is the particle
dictionary table.  Section~\ref{sec:qft} states the structural difference
from QFT.  Section~\ref{sec:dof} counts the degrees of freedom.
Section~\ref{sec:open} is an honest statement of open problems.

%=============================================================
\section{Four Excitation Types of $\Theta$}
\label{sec:types}
%=============================================================

\subsection{Type A — $\psi$-Taylor Modes (Imaginary-Time Expansion)}
\label{sec:typeA}

Because $\tau = t + i\psi$ and $\Theta$ is analytic in $\tau$, we may expand
around $\psi = 0$:
\begin{equation}
  \Theta(x, t, \psi) = \sum_{n=0}^{\infty} \frac{\psi^n}{n!}\,\Theta_n(x,t),
  \qquad
  \Theta_n := \frac{\partial^n\Theta}{\partial\psi^n}\bigg|_{\psi=0}.
  \label{eq:taylor}
\end{equation}
Each $\Theta_n$ is an ordinary $\B$-valued field on $M^4$.  The
identification is:

\begin{itemize}
  \item $n=0$: first fermionic generation (electron family) \conjecture
  \item $n=1$: second generation (muon family) \conjecture
  \item $n=2$: third generation (tau family) \conjecture
  \item $n\ge 3$: heavier, increasingly unstable modes;
        no stable fourth generation observed to date \conjecture
\end{itemize}

\paragraph{Physical reason for three stable generations.}
The $\psi$-circle is compact, $\psi\in[0,2\pi]$.  Factorial suppression in
\eqref{eq:taylor} causes higher-$n$ modes to carry progressively smaller
residues, explaining the observed mass hierarchy.  Furthermore, a discrete
$\psi$-\emph{parity} symmetry $\psi\mapsto -\psi$ is proved to be a symmetry
of the UBT action \proved\ (see \texttt{st3\_complex\_time\_generations.tex},
Sec.~3).  Under this parity $\Theta_n\mapsto(-1)^n\Theta_n$, so even-$n$
modes ($\Theta_0,\Theta_2$) and odd-$n$ modes ($\Theta_1$) belong to
different superselection sectors.  In particular $\Theta_0$ carries conserved
$\psi$-parity and cannot decay to lower modes \conjecture.  This provides a
topological stability argument for the electron analogous to charge
conservation.

\paragraph{Same quantum numbers.}
It is proved \proved\ that each mode $\Theta_n$ transforms identically under
any symmetry that acts only on the real sector $(M^4,t)$; in particular each
mode carries the same gauge quantum numbers under whatever gauge group emerges
from $\B$.  The mass difference between generations must therefore arise from
the $\psi$-sector dynamics, not from the gauge sector.

\paragraph{Quark generations.}
Quarks carry colour charge.  Whether quarks are also Type-A modes of $\Theta$
depends on whether colour ($\mathrm{SU}(3)_c$) is embedded in $\B$ or requires
octonionic extension.  This is \openq\ (see Section~\ref{sec:open}).

\subsection{Type B — Internal Algebra Projection}
\label{sec:typeB}

The biquaternion algebra $\B = \C\otimes_\R\H$ has 8 real basis elements,
which split under the Lorentz group into representations of definite spin:
\begin{equation}
  \{1\} \;\oplus\; \{i,j,k\} \;\oplus\; \{ij,ik,jk\} \;\oplus\; \{ijk\}
  \label{eq:alg_split}
\end{equation}
(using $i,j,k$ for quaternion imaginary units and $\overline{\cdot}$ for the
complex imaginary factor $\sqrt{-1}$).  The spin assignments are:

\begin{itemize}
  \item $\{1\}$ (scalar, spin 0) $\longrightarrow$ Higgs / scalar curvature
        \conjecture
  \item $\{i,j,k\}$ (vector, spin 1) $\longrightarrow$ photon, $W^\pm$, $Z$,
        gluons \conjecture\ (with status differing by particle; see
        Section~\ref{sec:dict})
  \item $\{ij,ik,jk\}$ (bivector, spin 1/2) $\longrightarrow$ fermions
        \conjecture
  \item $\{ijk\}$ (pseudoscalar/axial, spin 0 or 2?) $\longrightarrow$
        graviton candidate? \openq
\end{itemize}

\begin{remark}
The spin-$\tfrac{1}{2}$ character of the bivector sector follows from the
fact that the three imaginary quaternion units $\{i,j,k\}$ generate an
$\mathfrak{su}(2)$ algebra, and the bivectors $\{ij,ik,jk\}$ form the
$\mathbf{2}$ representation under the complexification of this algebra.
This is the algebraic source of spinors within $\B$. \conjecture
\end{remark}

Gauge bosons arise from phases on the $\psi$-circle:

\begin{itemize}
  \item $\mathrm{U}(1)_{\mathrm{EM}}$: identified with the global phase on the
        $\psi$-circle $e^{i\varphi}$, $\varphi\in[0,2\pi]$.  Localization
        of this phase yields the $\mathrm{U}(1)$ gauge field $A_\mu$.
        Maxwell's equations follow.  \proved\ (Appendix~C,
        \texttt{appendix\_C\_electromagnetism\_gauge\_consolidated.tex})
  \item $\mathrm{SU}(2)_L$: the left-action of unit quaternions on $\B$ is
        shown to be an $\mathrm{SU}(2)$ automorphism group.  Gauging this
        action gives $W^\pm$, $Z$.  \proved\ (Appendix~E2,
        \texttt{appendix\_E2\_SM\_geometry.tex}, Sec.~2)
  \item $\mathrm{SU}(2)_R\times\mathrm{U}(1)_Y$ from right-actions of
        unit quaternions: derived from $\B$ structure.  \proved\
  \item $\mathrm{SU}(3)_c$ (colour): \textbf{NOT} derived from $\B=\C\otimes\H$
        alone.  Requires octonionic extension $\C\otimes\mathbb{O}$ (Track~B).
        \openq
\end{itemize}

\subsection{Type C — Topological Configurations}
\label{sec:typeC}

When $\Theta$ is smooth on $M^4$ with boundary conditions
$\Theta\to\Theta_\infty$ as $|x|\to\infty$, space compactifies to $S^3$ and
$\Theta$ defines a map $S^3\to\text{target}$.  Different homotopy classes of
such maps are \emph{topological sectors}, stable against perturbative decay.

\begin{itemize}
  \item Winding number $n=0$: vacuum. \proved
  \item Winding number $n=1$ (simplest soliton): lightest stable topological
        excitation.  Candidate for electron at a mass scale determined by UBT
        coupling constants.  Relationship to Type-A $n=0$ mode is unclear.
        \openq
  \item \textbf{Hopfion} ($Q_H=1$, Hopf fibration $S^3\to S^2$): a knotted
        field configuration with conserved integer Hopf charge $Q_H\in\Z$.
        Stability is guaranteed by the Vakulenko--Kapitanski bound
        $E\ge c\,|Q_H|^{3/4}$ \proved\ (Appendix~I,
        \texttt{appendix\_I\_new\_fields\_and\_particles.tex}, Sec.~I.4).
        Proposed as candidate for proton / neutron: each nucleon is a
        $Q_H=\pm 1$ hopfion.  \openq\
        (See also \texttt{solution\_P5\_dark\_matter/dark\_matter\_hopfion\_solution.tex})
  \item Higher-charge hopfions $|Q_H|>1$: heavier topological objects.
        Candidate for excited baryons.  \speculative
\end{itemize}

\begin{remark}
Composite hadrons (proton, neutron) appear in UBT as topological defects, NOT
as $\psi$-Taylor modes.  This is structurally important: it means the proton
is \emph{not} a third-generation fermion but a qualitatively different object
— a soliton.  Proton stability then follows from topological charge
conservation, not from a gauge symmetry.
\end{remark}

\subsection{Type D — Real vs.\ Imaginary Projection}
\label{sec:typeD}

The biquaternion field has real and imaginary parts:
\begin{equation}
  \Theta = \re(\Theta) + i\,\im(\Theta),
  \qquad
  \re(\Theta), \im(\Theta) \in \H.
\end{equation}
In the $\psi\to 0$ limit, observable Standard Model fields arise from
$\re(\Theta)$.  The imaginary components $\im(\Theta)$ are present in the
algebra but decouple from classical SM observations:

\begin{itemize}
  \item $\re(\Theta)$: observable sector (SM particles, metric) \proved
  \item $\im(\Theta)$: dark sector.  These components contribute to the
        stress-energy tensor and hence to gravity but do not couple to the
        SM gauge fields.  Candidate for dark matter and dark energy.
        \speculative\
        (Appendix~U, \texttt{appendix\_U\_dark\_matter\_unified\_padic.tex})
\end{itemize}

%=============================================================
\section{The Particle Dictionary}
\label{sec:dict}
%=============================================================

Table~\ref{tab:dict} assigns each known particle (or particle class) to an
excitation type, a component of $\Theta$, and a proof status.  Where the
assignment is genuinely unknown, the table says so explicitly.

\begin{table}[htbp]
\centering
\caption{UBT Particle Dictionary. Status: \textbf{P}=Proved,
  \textbf{C}=Conjecture, \textbf{O}=Open, \textbf{S}=Speculative.}
\label{tab:dict}
\renewcommand{\arraystretch}{1.35}
\begin{tabular}{>{\raggedright}p{3.2cm}
                >{\raggedright}p{2.8cm}
                >{\raggedright}p{4.8cm}
                c}
\toprule
\textbf{Particle} & \textbf{Excitation type} & \textbf{$\Theta$ component} & \textbf{Status} \\
\midrule
Electron ($e^-$)
  & Type A, $n=0$
  & $\Theta_0$, $\psi$-even mode; bivector ($\{ij,ik,jk\}$) sector
  & C \\
Muon ($\mu^-$)
  & Type A, $n=1$
  & $\Theta_1$, $\psi$-odd mode; bivector sector
  & C \\
Tau ($\tau^-$)
  & Type A, $n=2$
  & $\Theta_2$, $\psi$-even mode; bivector sector
  & C \\
Electron neutrino ($\nu_e$)
  & Type A, $n=0$?
  & $\Theta_0$, neutral projection; exact component unknown
  & O \\
Muon neutrino ($\nu_\mu$)
  & Type A, $n=1$?
  & $\Theta_1$, neutral projection
  & O \\
Tau neutrino ($\nu_\tau$)
  & Type A, $n=2$?
  & $\Theta_2$, neutral projection
  & O \\
Up, charm, top quarks
  & Type A?, $n=0,1,2$?
  & $\Theta_{0,1,2}$ in colour sector; requires SU$(3)_c$
  & O \\
Down, strange, bottom quarks
  & Type A?, $n=0,1,2$?
  & $\Theta_{0,1,2}$ in colour sector; requires SU$(3)_c$
  & O \\
Photon ($\gamma$)
  & Type B, vector
  & U$(1)_{\mathrm{EM}}$ phase on $\psi$-circle; $\{i,j,k\}$ sector
  & P \\
$W^\pm$
  & Type B, vector
  & SU$(2)_L$ left-action generators on $\B$
  & P \\
$Z^0$
  & Type B, vector
  & SU$(2)_L \times$ U$(1)_Y$ mixing; derived from $\B$
  & P \\
Gluons ($g^a$, $a=1\ldots 8$)
  & Type B, vector
  & SU$(3)_c$ generators; \emph{not} from $\B$ alone
  & O \\
Graviton ($h_{\mu\nu}$)
  & Type B, tensor?
  & $\re(\mathcal{G}_{\mu\nu})$ fluctuation about vacuum metric
  & C \\
Proton ($p$)
  & Type C, topological
  & Hopfion $Q_H=+1$ configuration
  & O \\
Neutron ($n$)
  & Type C, topological
  & Hopfion $Q_H=-1$ or $Q_H=+1$ with isospin flip
  & O \\
Excited baryons
  & Type C, $|Q_H|>1$?
  & Higher-charge hopfions
  & S \\
Higgs boson ($H$)
  & Type B, scalar
  & Scalar part $\{1\}$ of $\Theta$; vacuum expectation drives SSB
  & O \\
Dark matter
  & Type D, imaginary
  & $\im(\Theta)$; decoupled from SM gauge sector
  & S \\
Dark energy
  & Type D, imaginary
  & Cosmological component of $\im(\Theta)$; p-adic sector
  & S \\
\bottomrule
\end{tabular}
\end{table}

\begin{remark}
Quark assignments are marked \openq\ because $\mathrm{SU}(3)_c$ colour
symmetry has not been derived from $\B = \C\otimes\H$ alone.  Quarks require
colour and therefore their Type-A identification is contingent on the
octonionic extension (Track B).  Treating quark generation indices as
$\psi$-Taylor modes is a plausible extrapolation of the lepton result but is
not independently established.
\end{remark}

%=============================================================
\section{Structural Difference from QFT}
\label{sec:qft}
%=============================================================

\subsection{Particle spectrum in QFT}

In quantum field theory the particle spectrum is defined by:
\begin{quote}
\emph{Particle spectrum = irreducible representations of the Poincaré group
$\times$ the gauge group $G_{\mathrm{SM}} = \mathrm{SU}(3)_c\times
\mathrm{SU}(2)_L\times\mathrm{U}(1)_Y$.}
\end{quote}
Each particle species corresponds to a separate field, and that field
transforms in a specific representation.  The gauge group and its
representations are inputs — they are \emph{postulated}, not derived.

\subsection{Particle spectrum in UBT}

In UBT the situation is structurally different:
\begin{quote}
\emph{Particle spectrum = excitation modes of the single field
$\Theta(x,\tau)\in\B$, where $\B = \C\otimes_\R\H$ with complex time
$\tau=t+i\psi$.}
\end{quote}
The Poincaré and gauge structure should \emph{emerge} from the geometry of
$\B$ and the complex time cylinder $\tau\in\C$, rather than being postulated
independently.

\paragraph{What has been derived (not postulated) in UBT.}
\begin{itemize}
  \item Lorentz group $\mathrm{SO}(3,1)$ from real quaternion rotations.
        \proved
  \item $\mathrm{U}(1)_{\mathrm{EM}}$ from phase on $\psi$-circle.  \proved
  \item $\mathrm{SU}(2)_L$ from left-unit-quaternion action on $\B$.  \proved
  \item $\mathrm{SU}(2)_R\times\mathrm{U}(1)_Y$ from right action.  \proved
  \item Three-generation structure via $\psi$-Taylor modes.  \conjecture
  \item Einstein's equations $G_{\mu\nu}=8\pi G T_{\mu\nu}$ as the real
        projection of the biquaternionic field equation.  \proved\
        (Appendix~R, \texttt{appendix\_R\_GR\_equivalence.tex})
\end{itemize}

\paragraph{What still requires postulate or extension.}
\begin{itemize}
  \item $\mathrm{SU}(3)_c$ (colour).  Requires $\C\otimes\mathbb{O}$.
        \openq
  \item Higgs mechanism (spontaneous symmetry breaking of electroweak).
        \openq
  \item Quark mass spectrum and CKM mixing.  \openq
  \item Neutrino masses and PMNS mixing.  \openq
\end{itemize}

\paragraph{Key structural contrast.}
In QFT: $\psi_e, A_\mu, q, \ldots$ are independent fields with independent
degrees of freedom.  The index structure of each is determined by the
representation theory of the postulated gauge group.

In UBT: $\Theta$ is a single object.  The different ``particles'' are
projections of $\Theta$ onto the different subalgebras, Taylor modes, and
topological sectors of one geometric object.  The independence of SM fields
is approximate: it holds to the extent that different excitation types
decouple in the low-energy limit.

%=============================================================
\section{Degree-of-Freedom Counting}
\label{sec:dof}
%=============================================================

\subsection{Excitation modes of $\Theta$}

\paragraph{Internal algebra (Type B).}
$\B = \C\otimes_\R\H$ has 8 real dimensions.  The basic decomposition
\eqref{eq:alg_split} gives $1+3+3+1 = 8$ basis elements.  With two
polarisations per spin-1 mode: $1\,(\text{scalar}) + 2\times 3\,
(\text{transverse vector}) + 4\,(\text{spinor}) + 1\,
(\text{pseudoscalar}) = 12$ physical modes per point.

\paragraph{$\psi$-Taylor modes (Type A).}
Infinite tower, but only the first $\sim 3$ are relevant due to factorial
suppression.  Each mode is an independent $\B$-valued field: $3\times 8 = 24$
internal d.o.f.\ for generations.

\paragraph{Extended matrix representation.}
The canonical extended field
$\Theta_{\mathrm{SM}}\in\C^{8\times 8}$ (\texttt{canonical/fields/theta\_field.tex})
has 128 real degrees of freedom.

\subsection{SM particle content}

A rough d.o.f.\ count for on-shell SM particles:
\begin{align}
  \text{6 quarks} \times 3\,\text{colours} \times 2\,\text{spins}
    &= 36\;\text{(plus antiparticles: 72)} \notag \\
  \text{6 leptons} \times 2\,\text{spins}
    &= 12\;\text{(plus: 24)} \notag \\
  \text{Gauge bosons: } \gamma(2)+W^\pm(6)+Z(3)+8g(16)
    &= 27 \notag \\
  \text{Higgs: }
    &= 4 \notag \\
  \text{Total (including antiparticles)}
    &\approx 127\;\text{physical d.o.f.}
  \label{eq:SMcount}
\end{align}

\subsection{Comparison}

The extended field $\Theta_{\mathrm{SM}}\in\C^{8\times 8}$ with 128 real
d.o.f.\ matches the SM count of $\approx 127$ to within one unit.  This
agreement is encouraging but should be treated as \conjecture: the mapping
between specific matrix entries and specific SM particles has not been
established in full generality.

The minimal field $\Theta\in\C^{4\times 4}$ (32 real d.o.f.) is sufficient
for one generation of SM particles without colour, consistent with the fact
that $\B = \C\otimes\H$ gives $\mathrm{SU}(2)\times\mathrm{U}(1)$ but not
$\mathrm{SU}(3)_c$.  Three generations require $3\times 32 = 96$ d.o.f.,
still below 128.  The remaining $128-96 = 32$ d.o.f.\ may accommodate gauge
bosons and the Higgs.

\begin{remark}
This counting does not prove that the identification is correct; it shows
that the $\C^{8\times 8}$ representation is at least \emph{large enough} to
contain the SM content.  A precise bijection between matrix entries and
particle quantum numbers remains \openq.
\end{remark}

%=============================================================
\section{Open Problems: An Honest Statement}
\label{sec:open}
%=============================================================

\subsection{What is proved}

\begin{enumerate}[label=(\alph*)]
  \item $\mathrm{U}(1)_{\mathrm{EM}}$ as a phase on the $\psi$-circle,
        with Maxwell's equations as a consequence.  (Appendix~C)
  \item $\mathrm{SU}(2)_L$ from left-action automorphisms of $\B$.
        (Appendix~E2)
  \item $\mathrm{SU}(2)_R\times\mathrm{U}(1)_Y$ from right action.
        (Appendix~E2)
  \item General Relativity as the real projection of the biquaternionic
        field equation.  (Appendix~R)
  \item $\psi$-parity is a symmetry of the UBT action; it places even and
        odd $\psi$-Taylor modes in different superselection sectors.
        (ST-3)
  \item The Vakulenko--Kapitanski bound $E\ge c\,|Q_H|^{3/4}$ ensures
        topological stability of hopfions.  (Appendix~I)
\end{enumerate}

\subsection{What is conjectured}

\begin{enumerate}[label=(\alph*)]
  \item Lepton generations $e,\mu,\tau$ correspond to $\psi$-Taylor modes
        $\Theta_0,\Theta_1,\Theta_2$.  (ST-3)
  \item The mass hierarchy $m_e < m_\mu < m_\tau$ arises from $\psi$-mode
        ordering with factorial suppression.
  \item The bivector sector of $\B$ provides the algebraic source of spin-$\tfrac{1}{2}$
        for fermions.
  \item The graviton is a linearised fluctuation of $\re(\mathcal{G}_{\mu\nu})$.
  \item The d.o.f.\ count of $\Theta_{\mathrm{SM}}\in\C^{8\times 8}$ ($128$
        real d.o.f.) matches the SM particle content.
\end{enumerate}

\subsection{What is open}

\begin{enumerate}[label=(\alph*)]
  \item $\mathrm{SU}(3)_c$: it does NOT follow from $\B = \C\otimes\H$ alone.
        Deriving colour requires the octonionic extension
        $\C\otimes\mathbb{O}$ (Track B in DERIVATION\_INDEX).  Whether this
        extension is natural or ad hoc is a central open question.
  \item Quark assignments: the identification of quarks as Type-A
        $\psi$-Taylor modes is a plausible extrapolation from leptons, but
        is not independently established; it is contingent on point (a).
  \item Neutrino masses: the PMNS mixing matrix and absolute neutrino mass
        scale have not been derived from UBT.
  \item Higgs mechanism: spontaneous symmetry breaking of
        $\mathrm{SU}(2)_L\times\mathrm{U}(1)_Y\to\mathrm{U}(1)_{\mathrm{EM}}$
        requires a scalar sector vacuum expectation value.  The scalar
        component of $\Theta$ is a candidate, but the full mechanism is
        not derived.
  \item Hadron structure beyond proton/neutron: the identification of
        higher baryons and mesons as higher topological sectors is
        speculative; no quantitative predictions exist yet.
  \item Precise bijection between $\Theta_{\mathrm{SM}}\in\C^{8\times 8}$
        matrix entries and SM particle quantum numbers.
\end{enumerate}

\subsection{What is unknown (fundamental)}

\begin{enumerate}[label=(\alph*)]
  \item Whether $\B = \C\otimes\H$ alone (without octonions) is sufficient
        for a complete description of nature, or whether $\C\otimes\mathbb{O}$
        is required.
  \item Whether UBT is renormalizable as a quantum field theory.
  \item Whether topological solitons (hopfions) can reproduce the proton
        mass from first principles ($m_p \approx 938\,\mathrm{MeV}$).
\end{enumerate}

%=============================================================
\section*{Summary}
%=============================================================

\noindent\textbf{The central claim} of UBT is that a single biquaternionic
field $\Theta$ with complex time gives rise to all observed particles through
four distinct excitation mechanisms.

\noindent\textbf{Strongest results:} $\mathrm{U}(1)_{\mathrm{EM}}$,
$\mathrm{SU}(2)_L$, GR — all derived, not postulated.

\noindent\textbf{Main conjecture:} three lepton generations from three
$\psi$-Taylor modes — structurally motivated, not yet a quantitative
prediction.

\noindent\textbf{Biggest gap:} $\mathrm{SU}(3)_c$ colour requires extension
beyond $\C\otimes\H$.

\noindent\textbf{This map turns conjectures into targets.} If
$\Theta_0\leftrightarrow e^-$ is correct, electron properties follow from
the $n=0$ mode geometry.  If proton~$\leftrightarrow$~hopfion is correct,
proton stability is topological.  If photon~$\leftrightarrow$~$\mathrm{U}(1)_{\mathrm{EM}}$
phase is correct (and it \emph{is} proved), Maxwell's equations follow from
UBT.

%=============================================================
\begin{thebibliography}{99}
\bibitem{appendix_C}
  D.~Jaro\v{s}, \textit{Appendix C: Electromagnetism and Gauge Structure},
  \texttt{consolidation\_project/appendix\_C\_electromagnetism\_gauge\_consolidated.tex}

\bibitem{appendix_E2}
  D.~Jaro\v{s}, \textit{Appendix E2: Standard Model Gauge Group from
  Biquaternionic Geometry},
  \texttt{consolidation\_project/appendix\_E2\_SM\_geometry.tex}

\bibitem{appendix_I}
  D.~Jaro\v{s}, \textit{Appendix I: Hopfions and Topological Field
  Configurations},
  \texttt{consolidation\_project/appendix\_I\_new\_fields\_and\_particles.tex}

\bibitem{appendix_R}
  D.~Jaro\v{s}, \textit{Appendix R: Recovery of General Relativity},
  \texttt{consolidation\_project/appendix\_R\_GR\_equivalence.tex}

\bibitem{appendix_U}
  D.~Jaro\v{s}, \textit{Appendix U: Dark Matter from p-adic and Imaginary
  Sector},
  \texttt{consolidation\_project/appendix\_U\_dark\_matter\_unified\_padic.tex}

\bibitem{st3}
  D.~Jaro\v{s}, \textit{ST-3: Three Generations from Complex Time},
  \texttt{research\_tracks/three\_generations/st3\_complex\_time\_generations.tex}

\bibitem{theta_field_md}
  D.~Jaro\v{s}, \textit{Formal Definition of the Biquaternionic Field
  $\Theta(q,\tau)$},
  \texttt{THETA\_FIELD\_DEFINITION.md}

\bibitem{canonical_theta}
  D.~Jaro\v{s}, \textit{Canonical Theta Field Definition},
  \texttt{canonical/fields/theta\_field.tex}
\end{thebibliography}

\end{document}
